% Options for packages loaded elsewhere
\PassOptionsToPackage{unicode}{hyperref}
\PassOptionsToPackage{hyphens}{url}
\documentclass[
  ignorenonframetext,
]{beamer}
\newif\ifbibliography
\usepackage{pgfpages}
\setbeamertemplate{caption}[numbered]
\setbeamertemplate{caption label separator}{: }
\setbeamercolor{caption name}{fg=normal text.fg}
\beamertemplatenavigationsymbolsempty
% remove section numbering
\setbeamertemplate{part page}{
  \centering
  \begin{beamercolorbox}[sep=16pt,center]{part title}
    \usebeamerfont{part title}\insertpart\par
  \end{beamercolorbox}
}
\setbeamertemplate{section page}{
  \centering
  \begin{beamercolorbox}[sep=12pt,center]{section title}
    \usebeamerfont{section title}\insertsection\par
  \end{beamercolorbox}
}
\setbeamertemplate{subsection page}{
  \centering
  \begin{beamercolorbox}[sep=8pt,center]{subsection title}
    \usebeamerfont{subsection title}\insertsubsection\par
  \end{beamercolorbox}
}
% Prevent slide breaks in the middle of a paragraph
\widowpenalties 1 10000
\raggedbottom
\AtBeginPart{
  \frame{\partpage}
}
\AtBeginSection{
  \ifbibliography
  \else
    \frame{\sectionpage}
  \fi
}
\AtBeginSubsection{
  \frame{\subsectionpage}
}
\usepackage{iftex}
\ifPDFTeX
  \usepackage[T1]{fontenc}
  \usepackage[utf8]{inputenc}
  \usepackage{textcomp} % provide euro and other symbols
\else % if luatex or xetex
  \usepackage{unicode-math} % this also loads fontspec
  \defaultfontfeatures{Scale=MatchLowercase}
  \defaultfontfeatures[\rmfamily]{Ligatures=TeX,Scale=1}
\fi
\usepackage{lmodern}
\ifPDFTeX\else
  % xetex/luatex font selection
\fi
% Use upquote if available, for straight quotes in verbatim environments
\IfFileExists{upquote.sty}{\usepackage{upquote}}{}
\IfFileExists{microtype.sty}{% use microtype if available
  \usepackage[]{microtype}
  \UseMicrotypeSet[protrusion]{basicmath} % disable protrusion for tt fonts
}{}
\makeatletter
\@ifundefined{KOMAClassName}{% if non-KOMA class
  \IfFileExists{parskip.sty}{%
    \usepackage{parskip}
  }{% else
    \setlength{\parindent}{0pt}
    \setlength{\parskip}{6pt plus 2pt minus 1pt}}
}{% if KOMA class
  \KOMAoptions{parskip=half}}
\makeatother
\setlength{\emergencystretch}{3em} % prevent overfull lines
\providecommand{\tightlist}{%
  \setlength{\itemsep}{0pt}\setlength{\parskip}{0pt}}
% Slide layout defaults for warning-clean scientific decks.
\usepackage{etoolbox}
\IfFileExists{xurl.sty}{\usepackage{xurl}}{}
\IfFileExists{seqsplit.sty}{\usepackage{seqsplit}}{\newcommand{\seqsplit}[1]{#1}}
\protected\def\breakseq#1{\seqsplit{#1}}
\protected\def\breaktt#1{\begingroup\ttfamily\seqsplit{#1}\endgroup}
\setlength{\emergencystretch}{6em}
\tolerance=5000
\hbadness=10000
\hfuzz=1pt
\setlength{\tabcolsep}{2pt}
\AtBeginEnvironment{longtable}{\tiny\renewcommand{\arraystretch}{0.86}\setlength{\tabcolsep}{1pt}}
\AtBeginEnvironment{tabular}{\tiny\renewcommand{\arraystretch}{0.86}\setlength{\tabcolsep}{1pt}}

% Natbib and cross-reference fallbacks — slides don't load natbib
% or cleveref, but combined-PDF manuscript prose may emit these
% commands. The fallback renders citations as a bracketed key list
% and cross-references as detokenized labels so slides don't fail on
% undefined control sequences or raw underscores. \providecommand is
% a no-op if packages load later.
\providecommand{\citep}[1]{[#1]}
\providecommand{\citet}[1]{#1}
\providecommand{\citealp}[1]{#1}
\providecommand{\citeauthor}[1]{#1}
\providecommand{\citeyear}[1]{#1}
\providecommand{\cref}[1]{\texttt{\detokenize{#1}}}
\providecommand{\Cref}[1]{\texttt{\detokenize{#1}}}

\newtheorem{proposition}{Proposition}
\newtheorem{hypothesis}{Hypothesis}
\usepackage{bookmark}
\IfFileExists{xurl.sty}{\usepackage{xurl}}{} % add URL line breaks if available
\urlstyle{same}
% fallback for those not using the hyperref driver hyperxmp:
\makeatletter
\@ifundefined{xmpquote}{\newcommand{\xmpquote}[1]{#1}}{}
\makeatother
\hypersetup{
  hidelinks,
  pdfcreator={LaTeX via pandoc}}

\author{\texorpdfstring{}{}}
\date{}

\begin{document}

\begin{frame}[fragile,allowframebreaks]{Stage 3: Nanopublication-Based
Knowledge Graph \label{sec:methods_kg}}
\protect\phantomsection\label{stage-3-nanopublication-based-knowledge-graph}
Stage 3 is the methodological core of this work: it transforms
unstructured abstracts into a structured, RDF-compatible knowledge graph
of scientific evidence. The stage encompasses four tightly coupled
operations: LLM-based assertion extraction, nanopublication packaging,
knowledge graph construction, and citation-weighted hypothesis scoring.

\begin{block}{LLM-Based Assertion Extraction}
\protect\phantomsection\label{llm-based-assertion-extraction}
We extract assertions by prompting a locally hosted LLM (Ollama
\citep{ollama2024}) to assess each paper's abstract against eight
standard hypotheses. The model receives a structured prompt containing
the paper title, abstract, and hypothesis definitions, and returns a
JSON array where each element specifies a hypothesis ID, direction
(supports, contradicts, neutral, or irrelevant), a confidence score
\(c \in [0, 1]\), and a reasoning string. Assertions marked
``irrelevant'' are discarded; confidence values are clamped to
\([0, 1]\); and responses are validated against the known hypothesis ID
set. Papers lacking abstracts are skipped. Detailed prompt engineering,
error taxonomy, and validation methodology are documented in the
\hyperref[sec:extraction_pipeline]{extraction pipeline section}.
\end{block}

\begin{block}{Nanopublication Schema and RDF Structure}
\protect\phantomsection\label{nanopublication-schema-and-rdf-structure}
Each assertion is encoded as a \textbf{nanopublication}
\citep{groth2010anatomy, kuhn2016decentralized}---a minimal,
self-contained, machine-readable unit of scientific evidence. Formally,
each nanopublication is a tuple \((p, h, d, c)\) where \(p\) is the
paper identifier, \(h\) the hypothesis identifier,
\(d \in \{\text{supports}, \text{contradicts}, \text{neutral}\}\) the
direction, and \(c\) the confidence. Provenance metadata records the LLM
model, UTC timestamp, and paper identifier.

The pipeline serializes nanopublications in two complementary formats:

\begin{enumerate}
\item
  \textbf{JSON Lines} (one JSON object per line) for efficient
  incremental checkpointing. Assertions are saved at configurable
  intervals (default: every 50 papers), enabling the pipeline to resume
  from where it left off after interruption without re-processing
  already-analyzed papers. Deduplication uses the composite key
  \((paper\_id, hypothesis\_id)\); re-runs with improved models
  overwrite stale results.
\item
  \textbf{RDF/TriG} per the nanopublication standard
  (\href{https://nanopub.net/}{nanopub.net}), producing four named
  graphs per nanopublication:
\end{enumerate}

\begin{table}[htbp]
\centering
\caption{RDF/TriG nanopublication structure. Each nanopublication contains four named graphs encoding the assertion, its provenance, and publication metadata per the nanopublication standard (\texttt{nanopub.net}).}
\label{tab:nanopub_schema}
\begin{tabular}{lll}
\toprule
\textbf{Named Graph} & \textbf{Content} & \textbf{Key Predicates} \\
\midrule
Head & Links the nanopub resource to its three component graphs & \texttt{np:hasAssertion}, \breaktt{np:hasProvenance}, \breaktt{np:hasPublicationInfo} \\
Assertion & The core scientific claim & \texttt{aif:asserts}, \texttt{aif:supports}/\texttt{aif:contradicts}, \texttt{aif:claim}, \texttt{aif:confidence}, \breaktt{aif:citationCount} \\
Provenance & How the assertion was generated & \breaktt{prov:wasGeneratedBy}, \breaktt{prov:generatedAtTime}, \breaktt{prov:wasAttributedTo}, \breaktt{prov:hadPrimarySource} \\
Publication Info & Metadata about the nanopublication itself & \texttt{dc:created}, \texttt{dc:creator}, \texttt{dc:license} \\
\bottomrule
\end{tabular}
\end{table}

The namespace \breaktt{http://activeinference.institute/ontology/}
(prefix \texttt{aif:}) defines all domain predicates; the
nanopublication schema (\breaktt{http://www.nanopub.org/nschema\#},
prefix \texttt{np:}) provides structural predicates; provenance uses
PROV-O (\breaktt{http://www.w3.org/ns/prov\#}); and Dublin Core
(\breaktt{http://purl.org/dc/terms/}) provides publication metadata. The
TriG output is suitable for publication to the decentralized
nanopublication network and aligns with FAIR data principles:
\textbf{F}indable via URI-based identification, \textbf{A}ccessible via
standard RDF protocols, \textbf{I}nteroperable through W3C-standard
serialization, and \textbf{R}eusable with explicit provenance and CC0
licensing.
\end{block}

\begin{block}{Knowledge Graph Construction}
\protect\phantomsection\label{knowledge-graph-construction}
The knowledge graph is an RDF-compatible directed graph with three node
types: \textbf{paper nodes} (metadata: title, abstract, authors, year,
venue, citation count, domain), \textbf{assertion nodes} (claim text,
direction, hypothesis ID, confidence), and \textbf{hypothesis nodes}
(the eight standard hypotheses). Edges encode five relations defined in
the schema:

\begin{itemize}
\tightlist
\item
  \texttt{aif:asserts} --- Paper \(\to\) Assertion
\item
  \texttt{aif:cites} --- Paper \(\to\) Paper
\item
  \texttt{aif:belongsTo} --- Paper \(\to\) Subfield
\item
  \texttt{aif:supports} --- Assertion \(\to\) Hypothesis
\item
  \texttt{aif:contradicts} --- Assertion \(\to\) Hypothesis
\end{itemize}

The graph is implemented with a dual backend: \texttt{rdflib}
\citep{rdflib2023} when available (preferred for semantic web
compatibility), with automatic fallback to \breaktt{networkx.DiGraph} for
environments without RDF dependencies. Both backends maintain identical
internal indices for efficient paper, assertion, and hypothesis queries.
\end{block}

\begin{block}{Citation-Weighted Hypothesis Scoring}
\protect\phantomsection\label{citation-weighted-hypothesis-scoring}
For each hypothesis \(H\), we compute a citation-weighted evidence
score:

\begin{equation}
\text{score}(H) = \frac{\sum_{a \in S(H)} w(a) - \sum_{a \in C(H)} w(a)}{\sum_{a \in A(H)} w(a)} \label{eq:score}
\end{equation}

where \(S(H)\), \(C(H)\), and \(A(H)\) are the sets of supporting,
contradicting, and all assertions for \(H\), and the weight function is:

\begin{equation}
w(a) = \log(1 + \text{citations}(a)) \cdot \text{confidence}(a) \label{eq:weight}
\end{equation}

The logarithmic citation weighting ensures that highly cited papers
carry more influence without allowing any single paper to dominate. The
score lies in \([-1, 1]\). \textbf{Interpretation note:} a score of
\(+0.7\) indicates that 70\% of weighted evidence supports the
hypothesis (net of contradictions and normalized by total weighted
evidence), \emph{not} that the hypothesis has a 70\% probability of
being true. Scores are best interpreted as relative rankings across
hypotheses and as temporal trajectories within a hypothesis, rather than
as absolute probability estimates. Temporal trends are computed by
evaluating the cumulative score at each year, using only assertions from
papers published up to that year. A full derivation appears in Appendix
\ref{sec:appendix_scoring}.
\end{block}

\begin{block}{Tally-Based Evidence Aggregation}
\protect\phantomsection\label{tally-based-evidence-aggregation}
We emphasize that this algorithmic scoring formula constitutes a
\textbf{tally-based approach} to evidence synthesis: each
nanopublication assertion operates as an independent evidential vote,
weighted by citation impact and the extraction model's confidence. The
aggregation is linear and additive---supporting and contradicting
assertions are summed and differenced without modeling dependencies,
correlated evidence, or causal structure among claims. This design
choice prioritizes transparency, reproducibility, and computational
tractability over statistical sophistication.

The tally-based framing introduces three constraints. First, assertions
from methodologically related papers (e.g., iterative publications from
a single research group testing the same model) are counted
independently, amplifying correlated evidence. To illustrate: if a group
publishes three papers (2019, 2021, 2023) reporting successively refined
variants of the same predictive coding model, each with high citation
counts, the scoring formula counts three independent supporting
assertions for H4---even though the underlying empirical evidence is
largely overlapping. An evidential diversity index (proposed in the
\hyperref[sec:conclusion]{conclusion}) would downweight this cluster.
Second, the scoring metric treats all assertion sources symmetrically:
an assertion from a theoretical review and one from an empirical trial
carry equal weight at a given confidence level. Third, temporal scoring
tracks \emph{cumulative totals} rather than dynamic probabilistic
estimates; the score at year \(t\) is the sum of all historical
evidence, rather than a decaying posterior that downweights early work.

We embrace these constraints intentionally. The tally-based approach
provides a stable, interpretable baseline against which more
sophisticated scoring methods can be evaluated. The
\hyperref[sec:conclusion]{conclusion} describes concrete
extensions---including hierarchical Bayesian scoring, causal evidence
graphs, and evidential diversity indices that downweight correlated
evidence.
\end{block}
\end{frame}

\end{document}
